% TOP_PROBLEM: {"id": 3900008, "title": "Integer-distance point sets", "category_id": 6, "category": {"id": 6, "name": "geometry", "display_name": "Geometry", "description": "Euclidean and non-Euclidean geometry, geometric structures.", "slug": "geometry", "order_index": 6, "created_at": "2026-07-31T15:26:25.671Z"}, "status": "open"}
% FIRST_POSTED: null
% ENTRY_KIND: statement_only
% Imported from: batch13.tex; UnsolvedMath ID 3900008
\documentclass[11pt]{book}
\usepackage{fontspec}
\setmainfont{Latin Modern Roman}
\setsansfont{Latin Modern Sans}
\usepackage{amsmath,amssymb,mathtools}
\usepackage{unicode-math}
\setmathfont{Latin Modern Math}
\usepackage{xurl}
\usepackage[hidelinks]{hyperref}
\newcommand{\sref}[2]{\hyperlink{source:#1:#2}{[S#2]}}
\newcommand{\cataloguescope}[1]{\paragraph{Mathematical question.}#1}
\begin{document}
% UnsolvedMath ID 3900008; source catalogue code AMR-038-0008
\setcounter{section}{3900007}
\section{Integer-distance point sets}\label{3900008}
{\small\sffamily UnsolvedMath ID 3900008 · Geometry}
\subsection{Definitions and mathematical statement}

\paragraph{Shared notation.}$\mathbb N=\{1,2,\ldots\}$, and $\mathbb Z,\mathbb Q,\mathbb R,\mathbb C$ denote the integers, rationals, reals and complex numbers. Any other conventions are specified locally.


Write a point of the Euclidean plane as $p=(x,y)\in\mathbb R^2$, with distance $|p-q|=\sqrt{(x-x')^2+(y-y')^2}$. A finite set $P$ is \emph{integral} if $|p-q|\in\mathbb N$ for all distinct $p,q\in P$. Here \emph{general position} means that no three points of $P$ lie on one affine line and no four lie on one Euclidean circle. It does not require integer coordinates or convex position.
For three distinct indexed points the noncollinearity condition is
\[
\det\begin{pmatrix}x_i&y_i&1\\x_j&y_j&1\\x_k&y_k&1\end{pmatrix}\ne0.
\]
Under this condition, four indexed points are nonconcyclic exactly when
\[
\det\begin{pmatrix}
x_i^2+y_i^2&x_i&y_i&1\\
x_j^2+y_j^2&x_j&y_j&1\\
x_k^2+y_k^2&x_k&y_k&1\\
x_\ell^2+y_\ell^2&x_\ell&y_\ell&1
\end{pmatrix}\ne0.
\]
The original paper supplies seven-point configurations with these properties, so the historically recorded existence question has an affirmative answer. \sref{3900008}{2}
\paragraph{Statement recorded in the supplied source.}
Are there seven distinct points $p_1,\ldots,p_7\in\mathbb R^2$ in general position such that $|p_i-p_j|\in\mathbb N$ for every $1\le i<j\le7$? \sref{3900008}{2}
\subsection{Short English statement}
Can seven points in the plane have all pairwise distances integral, with no three on a line and no four on a circle? Kreisel and Kurz constructed such sets; integer coordinates and convex position are not required.
\subsection{Sources}
\begin{description}
\item[\hypertarget{source:3900008:1}{S1}] ULAM.AI and the UnsolvedMath Contributors, \emph{UnsolvedMath: A Curated Collection of Open Mathematics Problems}, record 3900008 (AMR-038-0008). CC BY 4.0. \url{https://huggingface.co/datasets/ulamai/UnsolvedMath}.
\item[\hypertarget{source:3900008:2}{S2}] Tobias Kreisel and Sascha Kurz, \emph{There are integral heptagons, no three points on a line, no four on a circle}, Discrete \& Computational Geometry 39 (2008), 786--790, arXiv:0804.1303. Abstract; section~1, first paragraph; section~2, distance matrix~(1), Theorem~1 and its exact planar embedding (Figure~1). \url{https://arxiv.org/abs/0804.1303}.
\end{description}
\label{end:3900008}
\end{document}
